\section{Voice emerges from situated attention}
\label{sec:voice}

Of the distinctions in this article, situated attention is the one I find most
useful. It says more about the writing than a typo does. A writer notices the
detail, tension, or exception that matters to this subject and this reader,
then allocates space accordingly. The resulting voice depends on judgment
before it depends on ornament.

\subsection{Supported specifics resist generic paraphrase}

Names, measurements, dates, odd quotations, and concrete observations narrow
the range of plausible continuations. They also give the reader something to
verify. Agent-like prose can round such a detail into broad significance,
and an overzealous editor can repeat the loss while improving the flow.

Specificity does not license invention. A vivid anecdote, sensory detail, or
opinion requires the same support as a quantitative claim. The useful detail
belongs either to the writer's experience or to the available evidence.

\subsection{Unresolved tension can remain unresolved}

People often hold mixed views. A writer may consider a policy broadly useful
and still distrust one consequence; a study may show a strong result in one
regime alongside an unexplained failure in another. When the evidence does
not settle the tension, preserving it is more accurate than extracting a neat
lesson.

An agent-generated continuation may turn uncertainty into a balanced
conclusion, answer criticism with an optimistic final paragraph, or generalize
from an unresolved observation. Revision can end at the unresolved fact
because completeness of tone is not completeness of evidence.

\subsection{Deliberate irregularity carries information}

A distinctive word, abrupt transition, or parenthetical aside may look
inefficient in isolation. Its value depends on whether the writer can account
for the effect: perhaps it registers surprise, changes the level of
confidence, or lets a technical caveat interrupt a broader claim. Deliberate
irregularity carries voice when the reader can infer why it is there. Random
variation has no such account.

Era-bound references, local language, and field-specific phrasing can carry
the same function. Their intelligibility depends on the audience. An
explanation may be necessary for a broad readership, but replacement with
generic internet language usually removes more than it clarifies.

\subsection{Polish and provenance are different variables}

Perfect grammar can come from a human author or editor. Formal vocabulary, a
dry register, letter-style greetings, and common transitions are similarly
ambiguous. One dash or one short sentence says little, and unsourced prose is
a research problem rather than an authorship test. Templates account for much
of the regularity in correctly formatted documents.

Secondhand language requires separate treatment because a suspicious phrase
inside a quotation or title belongs to its source. Proper names and worked
examples have the same protection. Their wording can be explained or cited,
but revision would corrupt the evidence under discussion.

Version history supplies stronger evidence than stylistic guesswork. Text
committed before ChatGPT's public release on November 30, 2022, could not have
been produced with ChatGPT at that time. Drafts and revision records provide
similar provenance, which should outweigh a detector's impression of the final
surface.

\subsection{Genre sets the range of available personality}

Essays and personal writing can support humor, first person, and openly mixed
feelings. Technical or legal prose may require a neutral voice, yet neutrality
need not become generic. Humanization succeeds when it restores the situated
choices available within the genre rather than adding personality
indiscriminately.
